\begin{tikzpicture}[font=\small, >=Stealth, c/.style={draw, minimum size=6mm}]
\def\inp{{{1,0,2,1,0},{0,1,3,1,1},{2,1,0,0,2},{1,3,1,2,0},{0,1,2,1,1}}}
\def\ker{{{1,0,-1},{1,0,-1},{1,0,-1}}}
\foreach \r in {0,...,4} \foreach \c in {0,...,4} {
  \pgfmathsetmacro{\v}{\inp[\r][\c]}
  \pgfmathtruncatemacro{\hl}{(\r<3)&&(\c<3)}
  \ifnum\hl=1 \node[c, fill=blue!20] at (0.6*\c,-0.6*\r) {\pgfmathprintnumber{\v}};
  \else \node[c] at (0.6*\c,-0.6*\r) {\pgfmathprintnumber{\v}}; \fi
}
\node at (1.2,0.7) {输入特征图 $5\times5$};
\node at (3.6,-1.2) {\Large $*$};
\foreach \r in {0,...,2} \foreach \c in {0,...,2} {
  \pgfmathsetmacro{\v}{\ker[\r][\c]}
  \node[c, fill=orange!20] at (4.4+0.6*\c,-0.6-0.6*\r) {\pgfmathprintnumber{\v}};
}
\node at (5.0,0.7) {卷积核 $3\times3$};
\node at (6.6,-1.2) {\Large $=$};
\foreach \r in {0,...,2} \foreach \c in {0,...,2} \node[c] at (7.4+0.6*\c,-0.6-0.6*\r) {};
\node[c, fill=green!30] at (7.4,-0.6) {$0$};
\node at (8.0,0.7) {输出特征图 $3\times3$};
\node[font=\scriptsize, align=left, anchor=north west] at (-0.3,-3.2) {左上角 $3\times3$ 窗口与卷积核逐元素相乘再求和：$(1{+}0{+}2) - (2{+}1{+}0) = 0$；\\窗口每次向右或向下滑动 1 格（步长为 1），得到下一个输出值。};
\end{tikzpicture}
